% !TEX program = xelatex
% English Workbench 参赛应用方案（iCAN 2026）
% 编译：xelatex main.tex （连续两次以稳定交叉引用）
\documentclass[11pt,a4paper,fontset=none]{ctexart}

% ---------- 中文字体：统一使用宋体 ----------
\setCJKmainfont{Songti SC}
\setCJKsansfont{Songti SC}
\setCJKmonofont{Songti SC}
\setCJKfamilyfont{zhsong}{Songti SC}
\newcommand{\songti}{\CJKfamily{zhsong}}

\usepackage[a4paper,margin=2.4cm,top=2.6cm,bottom=2.6cm]{geometry}
\usepackage{booktabs}
\usepackage{tabularx}
\usepackage{array}
\usepackage{ragged2e}
\usepackage{enumitem}
\usepackage{graphicx}
\usepackage{float}
\usepackage[section]{placeins}
\usepackage{needspace}
\usepackage{xcolor}
\usepackage{tikz}
\usetikzlibrary{positioning,arrows.meta,calc,fit,backgrounds,shapes.geometric}
\usepackage[colorlinks=true,linkcolor=accent2,citecolor=accent2,urlcolor=accent2]{hyperref}

% ---------- 配色 ----------
\definecolor{ink}{HTML}{24313D}      % 正文深色
\definecolor{line}{HTML}{A9B6C2}     % 边框
\definecolor{fill}{HTML}{F3F6F9}     % 浅底
\definecolor{soft}{HTML}{E8EEF4}     % 更浅底
\definecolor{accent}{HTML}{C08A2E}   % 暖金强调
\definecolor{accent2}{HTML}{35617E}  % 钢蓝强调
\definecolor{muted}{HTML}{6B7A88}

% ---------- 表格列型 ----------
\newcolumntype{Y}{>{\RaggedRight\arraybackslash}X}
\newcolumntype{L}[1]{>{\RaggedRight\arraybackslash}p{#1}}

% ---------- 标题样式 ----------
\ctexset{
  section={
    format=\Large\bfseries\color{accent2},
    name={,、},
    number=\chinese{section},
    aftername=\hspace{0.25em},
    beforeskip=1.4ex plus .3ex,
    afterskip=0.9ex plus .2ex
  },
  subsection={
    format=\large\bfseries\color{ink},
    number=\arabic{section}.\arabic{subsection},
    aftername=\hspace{0.35em},
    beforeskip=1.0ex plus .2ex,
    afterskip=0.6ex plus .1ex
  }
}

% ---------- 常用样式 ----------
\setlist[itemize]{leftmargin=1.6em,itemsep=2pt,topsep=3pt}
\setlist[enumerate]{leftmargin=1.8em,itemsep=2pt,topsep=3pt}
\renewcommand{\arraystretch}{1.28}
\linespread{1.08}

\tikzset{
  box/.style={draw=line,fill=fill,rounded corners=2pt,minimum width=9cm,
              minimum height=0.95cm,inner sep=4pt,text=ink,align=center},
  side/.style={draw=line,fill=white,rounded corners=2pt,minimum width=2.8cm,
               minimum height=0.8cm,inner sep=3pt,text=ink,align=center},
  arr/.style={-{Stealth[length=2.2mm]},draw=accent2,thick},
  ext/.style={-{Stealth[length=2.2mm]},draw=accent,dashed,thick},
  fstep/.style={draw=line,fill=fill,rounded corners=2pt,minimum width=11.6cm,
               minimum height=0.72cm,inner sep=4pt,text=ink,align=left},
  fb/.style={-{Stealth[length=2.2mm]},draw=accent,dashed,thick},
  g/.style={draw=accent2,fill=soft,rounded corners=2pt,minimum width=2.6cm,
            minimum height=0.8cm,text=ink,align=center},
  g2/.style={draw=line,fill=fill,rounded corners=2pt,minimum width=6cm,
             minimum height=0.8cm,text=ink,align=center},
  g3/.style={draw=accent,fill=soft,rounded corners=2pt,minimum width=9cm,
             minimum height=0.85cm,text=ink,align=center},
  card/.style={draw=line,fill=fill,rounded corners=3pt,minimum width=3.3cm,
               minimum height=2.4cm,text=ink,align=center,inner sep=5pt},
  tag/.style={draw=line,fill=white,rounded corners=2pt,inner sep=3pt,
               text=muted,font=\scriptsize,align=center}
}

\begin{document}

% ===================== 标题区 =====================
\begin{center}
{\LARGE\bfseries\color{ink} English Workbench}\\[3pt]
{\large 面向初中英语教师的本地教学工作台与证据驱动智能体}\\[5pt]
{\normalsize\color{muted} 参赛应用方案}
\end{center}

\vspace{2pt}
\begin{center}
\setlength{\fboxsep}{8pt}%
\colorbox{soft}{\parbox{0.9\textwidth}{\centering\bfseries\color{ink}
项目主张：让数据替老师干活，让成长被学生看见。}}
\end{center}
\vspace{4pt}

\begin{center}
\begin{tabularx}{\textwidth}{@{}L{3.1cm}Y@{}}
\toprule
\textbf{作品名称} & English Workbench（英语教学工作台） \\
\textbf{AI 助手名称} & TeachMate \\
\textbf{参赛赛事及赛道} & iCAN 大学生创新创业大赛 · 华南赛区 · 高校组（软件作品，附软件著作权证明） \\
\textbf{学校、团队、成员} & 西安外国语大学；团队名称：TeachMate；成员：唐嘉钧（队长） \\
\textbf{当前版本} & 0.9.0-beta.2 \\
\textbf{作品形态} & 本地运行的软件与教学智能体应用 \\
\textbf{源码地址} & \url{https://github.com/JiajunTang2006/English-workbench} \\
\textbf{文档日期} & 2026 年 9 月 26 日 \\
\bottomrule
\end{tabularx}
\end{center}

\vspace{2pt}
{\small\tableofcontents}

\clearpage

% ===================== 项目概述 =====================
\section*{项目概述}
\addcontentsline{toc}{section}{项目概述}

初中英语教学既要参考日常学习表现，也要结合阶段性测评结果，教师需要据此持续判断学生掌握了什么、学得怎么样、有没有进展。可现实里，教学数据常常散落在表格、纸质材料和原卷附件中，统计结果很难直接转化成教学建议，过程性反馈也不够充分。本项目正围绕这些问题，探索一套面向单教师使用的教学辅助方案。

English Workbench 把教学数据管理、智能分析和成长反馈集中在同一个本地工作台。数据校验与统计计算交给业务程序完成；TeachMate 智能体则在教师指定的班级、考试、学生和资料范围内调用教学工具，产出带证据的分析报告与教学建议，再由教师审核后用于讲评、辅导和复习安排。成长森林用可追溯的活动记录，呈现学习参与、持续投入以及相对自身基线的进步，为过程性评价提供参考。

目前的 Beta 原型已实现班级与学生管理、成绩与过程性记录、原卷与错题管理，以及考试分析、学生诊断、复习计划和成长森林等功能，并完成了相应的工程检查。项目希望减轻重复性数据处理的负担、支持差异化教学、丰富成长反馈；教学应用效果会在后续试用中，通过任务完成时间、教师对建议的修改与采纳情况、复测记录等来评估。

% ===================== 一 =====================
\section{教育应用背景与项目目标}

\subsection{过程性评价对教学数据整合的需求}

英语教学中的学习表现有多种记录形式：默写反映一定范围内的词汇与拼写，背诵、写作和作业记录提供日常学习信息，考试则呈现特定时间与任务条件下的测评结果。这些信息的用途和解释边界各不相同。面向过程性评价的工具，需要保留每类记录的来源、时间和评价口径，并支持教师在学生、班级和学期之间建立关联。

一旦记录分散在表格、纸质材料和原卷附件里，教师就得额外做整理、核对和关联，才能拼出连续的学情判断。因此本项目把数据整合作为基础任务：以学期、班级和学生为共同关联，统一保存考试与日常记录，为后续分析提供一致、可核查的数据来源。

\subsection{差异化教学对分析与反馈的需求}

平均分、分数段和排名能描述成绩分布，但要制定讲评和辅导方案，还得进一步看题目表现、知识点覆盖和学生之间的差异。同一个总分，得失分结构可能完全不同；只有拿到相应的题目和作答证据，教师才能判断复习重点。

所以项目关注的是从“统计描述”到“教学行动”的转化：把可计算的事实与有依据的解释区分开，由智能体整理出分析发现和建议，同时保留教师审核与调整的环节。在这一过程中，AI 承担数据解释、方案草拟等辅助工作；建议是否适用，仍要结合教学目标、课时安排和学生实际情况来判断。

\subsection{成长反馈对持续学习支持的需求}

阶段性成绩只是评价信息的一部分。学生有没有订正、是否坚持复习、在可比任务里有没有进步，同样能为教师提供过程反馈的依据。为此，项目引入成长森林，把学习活动和相对自身基线的进步记录成可解释的成长事件，并与学业表现分开呈现。

这样做的目的是让反馈更具体：教师能说清学生完成了哪些学习活动、在哪些方面取得了进步、下一阶段还需要哪些支持。成长记录不作为学生综合能力的单一评分，它的激励效果和长期适用性有待后续试用检验。

\subsection{项目目标与实施范围}

围绕上述需求，项目设定了三个相互衔接的目标。

\textbf{数据管理目标：} 建立统一的教学记录组织方式，减少重复维护与人工汇总，为统计、分析和历史比较提供一致的数据基础。

\textbf{教学支持目标：} 把统计结果、题目材料和教学知识条目组织成可追溯的分析依据，辅助教师形成讲评、个别辅导和复习安排。

\textbf{成长反馈目标：} 记录学习参与、持续投入和可比进步，支持教师给出具体的过程反馈，同时保留奖励来源和更正依据。

项目当前以单教师本地使用为实施范围，优先适配已有的教学流程。功能设计服务于教学效率与反馈质量的改善，并不以增加练习数量作为成效指标。

% ===================== 二 =====================
\section{教学问题、目标用户与需求分析}

\subsection{教学流程中的主要问题}

\begin{table}[htbp]
\centering
\small
\begin{tabularx}{\textwidth}{@{}L{3.1cm}YY@{}}
\toprule
\textbf{痛点} & \textbf{对教师和学生的影响} & \textbf{本项目的应对方式} \\
\midrule
数据分散、重复录入 & 同一学生的信息需要在多处维护，容易遗漏或混淆 & 以学期、班级、学生和考试为共同关联，集中保存记录 \\
统计与解释脱节 & 知道平均分变化，却难以形成有依据的讲评安排 & 把统计事实提供给教学智能体，生成可核查的发现和建议 \\
讲评缺少针对性 & 不同学生接受相同的讲解和练习，教师难以兼顾差异 & 提供班级分析和单人、批量学生诊断 \\
错题难以继续利用 & 原卷、错题与后续训练之间联系薄弱 & 保留错题来源与考点，支持结合证据制定复习计划 \\
反馈过度依赖分数排名 & 部分学生的参与和进步不易被看到 & 用成长森林展示持续投入、自身进步及记录依据 \\
工具部署与数据管理有门槛 & 普通教师难以维护复杂服务，也担心数据失控 & 单教师本地部署，提供备份、恢复和明确的 AI 联网边界 \\
\bottomrule
\end{tabularx}
\caption{教学流程中的主要问题与本项目的应对方式}
\end{table}

上述问题构成了当前方案的设计依据；它们的适用范围和优先级，会在后续教师试用中进一步检验。

\subsection{用户范围与使用角色}

\textbf{核心使用者是初中英语教师。} 他们通常要带一个或多个班级，日常工作包括成绩整理、日常评价、讲评备课和个别辅导。系统围绕这些任务组织功能入口，让教师在既有流程里就能用上数据管理与智能分析能力。

\textbf{学生是教学反馈和成长支持的主要受益者。} 当前版本由教师操作，学生通过教师提供的画像反馈、复习安排和成长记录了解自己的学习表现；独立学生端属于后续规划。

\textbf{家长和学校是延伸受益者。} 教师可以在家校沟通中使用经过核对的分析材料，学校也可以把项目当作轻量教学工具试用。当前尚未建设独立的家长端、学校管理端或多人协同权限体系。

\subsection{功能需求与验收依据}

\begin{table}[htbp]
\centering
\small
\begin{tabularx}{\textwidth}{@{}L{3.1cm}YY@{}}
\toprule
\textbf{需求} & \textbf{用户期望} & \textbf{验收关注点} \\
\midrule
快速建档与录入 & 复用已有名单和成绩表，减少重复劳动 & 班级、姓名、学号、分数和缺考状态能正确对应；异常输入有提示 \\
一致的成绩统计 & 相同数据得到一致结果 & 分母、满分、阈值和缺考处理有明确规则 \\
有依据的 AI 分析 & 能知道结论依据什么数据 & 发现与建议引用当前任务的证据，数据不足时说明限制 \\
个性化教学反馈 & 同时看到优势、薄弱点和下一步 & 建议对应具体学生和考试，避免仅凭总分推断学习态度 \\
教师保留决定权 & 可以调整 AI 建议、控制任务范围 & 关键范围、批量任务和评价确认有明确操作入口 \\
成长可追溯 & 能解释学生为什么获得成长奖励 & 记录有来源，重复同步不重复奖励，更正可以追溯 \\
数据可保管、可恢复 & 换版本或误操作后仍能找回数据 & 具备数据库迁移、备份校验和恢复流程 \\
\bottomrule
\end{tabularx}
\caption{功能需求与验收依据}
\end{table}

\subsection{典型应用场景}

\textbf{考试后的讲评准备。} 教师导入成绩，确认满分、分层线和缺考情况；先看基础统计，再让 TeachMate 汇总班级薄弱点，结合原卷和逐题数据形成讲评建议。

\textbf{日常教学记录。} 教师记录默写、背诵、写作和作业情况，通过集中视图了解近期学习表现，需要时补充观察或订正记录。

\textbf{个别辅导与家校沟通。} 教师选定学生和考试，生成诊断草稿，核对其中的优势、问题和建议，再作为谈话或沟通材料使用。

\textbf{阶段复习。} 教师把分析结果转成复习安排，组织练习与复测，并把新的表现继续记录到工作台。目前的闭环仍包含教师组织实施和确认的环节。

% ===================== 三 =====================
\section{系统功能设计与实现情况}

\subsection{教学数据工作台}

\textbf{班级、学期与学生管理。} 教师按学期维护班级和学生档案，并查看学生的历史表现。数据按学期组织，减少新学期设置覆盖历史记录的风险。

\textbf{成绩管理与统计。} 支持录入和导入考试成绩，维护满分、缺考状态、优秀及合格阈值，以及 A/B/C 分层线等考试设置；平均分、分布、排名和相关统计由程序计算并展示。年级比较和逐题分析，分别以相应年级数据和逐题记录的完整性为前提。

\textbf{默写、背诵、写作与日常作业。} 按轮次或任务记录过程性表现，让教师把平时记录和考试结果放在同一工作空间查看。各类记录保留各自原有的含义——比如作业提交情况，并不等同于英语能力水平。

\textbf{原卷与错题。} 管理试卷附件和错题记录，保留题号、题型、考点、来源原卷及必要说明。附件与错题记录共同构成后续分析的数据来源；复杂的扫描件、手写答案和表格，识别结果仍需人工核验。

\textbf{仪表盘与待办。} 仪表盘提供当前教学数据的概览；待办事项记录需要继续跟进的事务，帮教师把数据整理和日常教学衔接起来。

\subsection{TeachMate 教学智能体}

\textbf{考试分析。} 根据当前考试的统计、题目和资料，整理班级表现、分数分布和有数据支撑的薄弱点，并形成教学建议。只有总分时，分析聚焦整体表现；具备逐题成绩、题目内容和考点时，才能进一步讨论具体题目和知识点。

\textbf{学生诊断。} 面向选定学生或一组学生，整合历史成绩与本次证据，形成学生画像和个性化建议。批量诊断会先生成执行方案，由教师检查范围与预计消耗后确认。诊断同时呈现优势与待改进之处，并把结论限定在现有记录能够支撑的范围内。

\textbf{复习计划。} 依据已有分析结果和教学知识条目，整理薄弱知识点的复习优先级、练习类型建议和时间安排，教师可结合课时和实际情况修改后使用。当前功能侧重复习安排与练习建议；自动组卷、题目质量审核和判分流程仍需完善。

\textbf{报告核对与教师确认。} 报告呈现分析发现、建议及其对应证据。教师可以核对原始记录、调整评价内容，再确认形成教学材料。系统区分 AI 原始草稿与教师确认内容，让生成、审核和使用各环节的责任边界保持清晰。

\subsection{英语成长森林}

成长森林用树木阶段和事件记录来呈现学习过程，为教师提供关于学习参与、自身进步和持续投入的反馈依据。

奖励由基础分、进步分和里程碑分组成，并受类别、每日和每周上限的约束。进步奖励依据学生自身的可比测评记录计算，规则并不直接按绝对分数排序；退步不触发扣分，缺失数据与有效零分也分开处理。

成长记录与学业表现分开解释：前者反映学习活动和投入，后者需要相应的测评证据。当前的学业表现分为词汇、语法、阅读、听力和写作五个分支，这属于本项目的产品分类，而非课程标准规定的核心素养分类。各分支遵循“证据与评价对象相对应”的原则，例如默写记录只支持相关词汇或拼写表现的判断。

教师可以查看奖励来源，并进行有记录的更正。考试归档、恢复和出勤状态变化，系统都会同步处理对应奖励；重复同步应保持结果一致。需要说明的是，成长森林是过程反馈工具，它对学习动机的实际影响仍需通过试用观察。

\begin{figure}[htbp]
\centering
\begin{tikzpicture}[font=\footnotesize,node distance=0.55cm and 0.8cm]
\node[g] (b) at (-3.4,0) {基础分};
\node[g] (p) at (0,0) {进步分};
\node[g] (m) at (3.4,0) {里程碑分};
\node[g2] (sum) at (0,-1.2) {合计 $r = $ 基础分 $+$ 进步分 $+$ 里程碑分};
\node[g2] (cap) at (0,-2.25) {封顶：类别上限 · 每日上限 · 每周上限};
\node[g3] (nut) at (0,-3.3) {学期营养 $=\max(0,\ $有效奖励 $+$ 更正净调整$)$};
\draw[arr] (b) -- (sum);
\draw[arr] (p) -- (sum);
\draw[arr] (m) -- (sum);
\draw[arr] (sum) -- (cap);
\draw[arr] (cap) -- (nut);
\node[tag] at (-4.6,-2.25) {退步不扣分};
\node[tag] at (4.6,-2.25) {缺失不补零};
\node[tag] at (0,-4.15) {基础分人人相同；进步分相对学生自身可比得分率；里程碑分奖励连续达成周数};
\end{tikzpicture}
\caption{成长森林的奖励构成与封顶逻辑（计分公式只在后端规则模块中实现一份）}
\end{figure}

\needspace{15\baselineskip}
\subsection{功能状态说明}

\begin{table}[htbp]
\centering
\small
\begin{tabularx}{\textwidth}{@{}L{4.6cm}L{2.2cm}Y@{}}
\toprule
\textbf{功能} & \textbf{当前状态} & \textbf{使用条件或后续工作} \\
\midrule
学期、班级、学生、考试管理 & 已有实现 & 由教师确认导入数据和设置 \\
默写、背诵、写作、作业记录 & 已有实现 & 由教师录入、导入或维护 \\
成绩统计、原卷和错题管理 & 已有实现 & 分析深度受数据完整度影响 \\
AI 考试分析、学生诊断、复习计划 & 已有实现 & 需要模型配置和可用服务；建议需教师审核 \\
批量诊断和教师确认 & 已有实现 & 先确认学生范围、执行方案及预算 \\
成长森林及可追溯奖励 & 已有实现 & 规则初值需试点调整，用于过程反馈 \\
数据备份、恢复和迁移 & 已有实现 & 完整备份需同时保留数据库、附件等文件 \\
个性化补偿题、过关测试的完整自动闭环 & 后续完善方向 & 需补齐题目质量审核、实施、判分和复测验证 \\
交互音标手册、听力 MP3 生成 & 规划功能 & 尚未确认完整实现，列入后续规划 \\
学生端、家长端、学校多人协同 & 规划方向 & 需要独立产品设计和访问权限体系 \\
\bottomrule
\end{tabularx}
\caption{功能状态说明}
\end{table}

% ===================== 四 =====================
\section{AI 在作品中的核心作用}

\subsection{智能分析的任务定位}

系统把确定性的业务计算与模型生成任务分开组织：平均分、得分率和分数分布由业务程序计算；TeachMate 负责理解教师意图、选择授权工具、整合当前任务的事实与参考条目，并生成分析和建议。这样的分工，让统计来源与解释过程可以分别核查。

以“分析当前考试并制定下周复习安排”为例：系统先依据教师选定的考试和班级确定数据范围，再获取统计与题目证据，按需补充资料，最后形成与证据关联的复习建议。智能体的核心作用，就体现在任务组织、工具调用和证据整合上，并通过教师审核衔接到实际教学。

\subsection{智能体运行流程}

\begin{table}[htbp]
\centering
\small
\begin{tabularx}{\textwidth}{@{}L{2.4cm}YY@{}}
\toprule
\textbf{步骤} & \textbf{智能体及系统行为} & \textbf{教师可以观察或控制的内容} \\
\midrule
1. 接收请求 & 接收自然语言问题及教师选择的上下文 & 当前学期、班级、考试、学生和已选资料 \\
2. 确定任务 & 进入考试分析、学生诊断或复习计划等能力流程 & 任务类型、分析范围；批量任务的执行方案 \\
3. 读取证据 & 调用受控工具获取统计、题目、学生成绩或正式附件 & 工具执行状态和相应证据 \\
4. 分析与补充 & 基于已有事实组织发现；必要时查询教学知识条目 & 处理中状态、所用资料和数据不足提示 \\
5. 检查输出 & 校验结构化结果和证据引用，处理不合格输出 & 报告、限制说明或明确的失败信息 \\
6. 教师确认 & 教师核对、修改评价，决定如何使用建议 & 最终采纳内容及后续教学安排 \\
\bottomrule
\end{tabularx}
\caption{智能体运行流程与教师的观察控制点}
\end{table}

\begin{figure}[htbp]
\centering
\begin{tikzpicture}[font=\footnotesize,node distance=0.34cm]
\node[fstep] (s1) {\textbf{1. 接收请求}\quad 接收自然语言问题及教师选择的上下文};
\node[fstep,below=of s1] (s2) {\textbf{2. 确定任务}\quad 进入考试分析 / 学生诊断 / 复习计划等能力流程};
\node[fstep,below=of s2] (s3) {\textbf{3. 读取证据}\quad 调用受控工具获取统计、题目、成绩或正式附件};
\node[fstep,below=of s3] (s4) {\textbf{4. 分析与补充}\quad 组织分析发现，必要时查询教学知识条目};
\node[fstep,below=of s4] (s5) {\textbf{5. 检查输出}\quad 校验结构化结果与证据引用，处理不合格输出};
\node[fstep,below=of s5] (s6) {\textbf{6. 教师确认}\quad 核对、修改评价，决定如何使用建议};
\draw[arr] (s1)--(s2);
\draw[arr] (s2)--(s3);
\draw[arr] (s3)--(s4);
\draw[arr] (s4)--(s5);
\draw[arr] (s5)--(s6);
\draw[fb] (s6.east) .. controls +(1.6,0) and +(1.6,0) .. node[right,text=accent,font=\scriptsize,align=left,pos=0.5]{教师可随时调整范围\\或要求重新生成} (s1.east);
\end{tikzpicture}
\caption{TeachMate 智能体的运行流程}
\end{figure}

系统把任务和工具事件记录成可观察的运行过程，前端通过进度与结果视图反馈状态。这些运行记录说明了请求处理、数据读取和结果生成之间的联系，为操作复核与异常定位提供依据。

\subsection{教学工具与知识支持}

教学桥接层提供考试分析数据包、考试概览、分数分布、题目统计、学生趋势、学生逐题成绩、正式附件读取、教学依据查询和报告提交等工具。数据范围由教师选择、由服务器确定，模型工具参数不能自行扩展到其他班级或学生。

知识支持采用内置教学条目与编译后的摘要、索引：系统优先给出当前任务的相关摘要，再按考点或条目补充具体依据，以控制输入规模。知识条目保留来源和可信度说明，并区分原始材料中的分类，与项目整理的考点映射、错因和干预建议。

报告发现和教学建议通过证据编号关联到当前任务的数据。例如“优先复习阅读定位”这一建议，需要对应本次读取的题目或成绩证据。引用校验负责检查证据是否存在及其归属，教师审核则进一步判断证据是否足以支撑相应解释——两者分别承担形式检查和教学判断的职责。

\subsection{教学归因与人工审核机制}

错题归因涉及词汇、语法、信息定位、审题和表达等因素，分数本身并不足以确定具体原因。因此，系统把错因分析定位为“基于现有证据的教学假设”，并通过限制说明呈现材料缺口。

只提供总分时，分析限于整体表现；有逐题得分时，可以讨论得失分结构；而涉及具体错因和学习过程的判断，还需要结合学生作答、课堂观察和交流信息。分析能走多深，取决于证据的完整程度。

教师负责确认分析范围、核对证据、修改建议并决定后续教学安排。系统的定位是教学辅助，目的是支持教师的专业判断与沟通。

% ===================== 五 =====================
\section{开发工具与技术实现方案}

\subsection{开发技术与职责}

\begin{table}[htbp]
\centering
\small
\begin{tabularx}{\textwidth}{@{}L{4.6cm}Y@{}}
\toprule
\textbf{技术或工具} & \textbf{在项目中的作用} \\
\midrule
Python、FastAPI、Uvicorn & 本地服务、教学业务接口和文件处理 \\
SQLite、SQLAlchemy、Alembic & 数据保存、业务数据访问及版本迁移 \\
HTML、CSS、JavaScript & 工作台界面、交互及图表相关页面逻辑 \\
SheetJS、openpyxl 等 & 表格导入、解析和导出相关处理 \\
PDF 文本及页面解析组件 & 为已上传资料提供文本和页面处理能力 \\
Node.js、DeepSeek Harness & 在配置完整时运行教学智能体及工具调用流程 \\
项目自建 Education Bridge & 将智能体调用限制到受控教学工具和当前数据范围 \\
模型接口适配层 & 连接已配置的文本或视觉模型服务 \\
pytest、Node 测试、Playwright & 业务回归、前端逻辑和浏览器流程检查 \\
Git、GitHub 及平台构建脚本 & 版本管理、源码交付和桌面构建支持 \\
\bottomrule
\end{tabularx}
\caption{主要开发技术与职责}
\end{table}

作品基于现有框架和第三方组件开发。项目的主要工作包括教学数据模型、教学工作台交互、教学智能体接入、证据与范围控制，以及成长反馈设计；基础模型和 Harness 运行时来自外部服务或第三方项目。

在降低使用门槛上，项目做了两层内置：一是智能体运行环境尽量做到开箱即用——智能体所依赖的基础指令和工具调用已在软件内预置，教师配置模型服务时通常只需填入 API Key 即可使用，不必自行编写命令或脚本；二是 TeachMate 的各教学插件已内置任务提示词（prompt），教师发起分析时不需要手工描述分析思路，从而降低了非技术背景教师使用 AI 的障碍。

\subsection{系统架构与数据流}

系统由教师交互界面、本地业务服务、数据存储、智能体运行层和模型服务组成，各部分围绕数据操作与分析任务协同运行。

教师界面提交数据操作或分析请求后，本地服务完成权限检查、数据校验和业务计算，并维护数据库及附件关联。分析任务由智能体运行层组织，通过 Education Bridge 在指定范围内读取教学数据，再与配置的模型服务交互；模型返回的结果经过结构与证据引用检查后，由本地服务保存并供界面展示。

当前实现由 FastAPI 提供 WorkBench / TeachMate 页面，Harness 作为后台智能体运行进程，源码启动器依据运行时文件的就绪状态选择执行路径。具体环境要求和启动操作见第六节。

在扩展性上，系统从设计之初就为外部数据源和上层智能体预留了对接空间：既可以通过学校提供的数据接口拉取成绩等业务数据进入工作台统一分析，也可以把工作台中的教学数据对接到 WorkBuddy 等其他智能体，用于更大范围的编排与协作，避免教学工具被局限在单一应用内。

\begin{figure}[htbp]
\centering
\begin{tikzpicture}[font=\footnotesize]
% 本地边界（先绘制，位于底层）
\draw[dashed,draw=accent2!65,rounded corners=4pt] (-4.7,1.12) rectangle (4.7,-4.72);
\node[anchor=south west,text=accent2,font=\scriptsize] at (-4.7,1.16) {本地运行（教学数据保存在本机）};
% 分层
\node[box] (ui) at (0,0) {\textbf{教师交互界面}\quad WorkBench / TeachMate 页面};
\node[box,below=0.42cm of ui] (svc) {\textbf{本地业务服务（FastAPI）}\quad 权限检查 · 数据校验 · 业务计算};
\node[box,below=0.42cm of svc] (db) {\textbf{数据存储}\quad SQLite（成绩 · 过程性记录 · 错题）与本地附件};
\node[box,below=0.42cm of db] (agent) {\textbf{智能体运行层}\quad Harness $+$ Education Bridge（受控教学工具）};
\node[box,below=0.55cm of agent,draw=accent] (llm) {\textbf{模型服务}\quad 已配置的文本 / 视觉模型（可选云端）};
\draw[arr] (ui) -- (svc);
\draw[arr] (svc) -- (db);
\draw[arr] (db) -- (agent);
\draw[arr] (agent) -- (llm);
% 外部对接
\node[side] (school) at (6.75,-2.1) {学校数据接口};
\node[side] (other) at (6.75,-3.65) {WorkBuddy 等其它智能体};
\draw[ext] (school.west) -- (db.east);
\draw[ext] (db.east) -- (other.west);
\node[tag,text=accent] at (6.75,-1.25) {可选对接};
\end{tikzpicture}
\caption{系统总体架构与数据流}
\end{figure}

\subsection{数据可靠性设计}

\textbf{应用与数据分离。} 教学数据独立存放在应用数据目录，程序升级通过受控迁移维护既有记录；附件采用本地文件与数据库元数据关联的方式管理。

\textbf{控制并发修改。} 工作区保存时用修订信息识别冲突，遇到旧版本数据提交会提示重新加载，减少无意覆盖。

\textbf{备份与恢复。} 升级前进行数据库检查和备份，完整备份包含数据库、附件及校验信息。恢复前先校验备份并保留现有数据副本，避免把“复制单个数据库文件”误当成完整备份。

\textbf{记录更正过程。} 成长奖励采用可追溯事件与更正机制，对归档、恢复、缺考和改分引起的变化进行同步；规则版本和来源记录有助于解释历史结果。

\subsection{数据保护与联网边界}

系统采用本地存储与可选云端推理相结合的部署方式。依赖安装完成后，日常记录和已有数据的统计都可以在本机运行；调用云端模型时，当前任务所需的上下文会发送到所配置的服务，因此仍需要网络，也可能产生调用费用。

系统通过身份字段过滤、匿名编号映射和附件文本处理，减少学号、电话及本地标识等信息进入模型请求。当前教师工作模式允许部分学生姓名参与分析，所以本地存储并不等于全部身份信息都不会外传。使用前需要确认资料范围和模型服务配置；竞赛展示使用合成数据。

本地接口使用启动时生成的访问 Token，教学工具按服务器确定的学期、班级、考试和学生范围读取数据。API 密钥不应写入源码、报告或录屏；第三方数据连接只在有明确授权并完成配置时才启用，它并不是参赛演示的必要条件。

在可控性上，本地部署本身已经为数据控制打下基础；同时软件保留了用户自主扩展的空间——教师或开发者可以自行安装插件，也可以通过智能体对软件能力进行改造和增强，让工具随使用场景演进，而不必等待官方版本更新。

% ===================== 六 =====================
\section{使用说明与交互流程}

\subsection{安装与首次启动}

本项目采用公开源码交付，评审可通过 GitHub 仓库获取指定版本，运行前需准备相应依赖和模型服务配置。

首次运行需要 Python 3.11 环境，按仓库的依赖声明安装 Python 组件；使用完整 Harness 路径时，还需准备相匹配的 Node.js、包管理器及运行时构建产物。完成依赖准备后，macOS 可使用 \texttt{start\_macos.command}，Windows 开发环境可使用 \texttt{start\_windows\_dev.bat}，也可在项目根目录运行 \texttt{python launcher.py}。

启动器会启动本地服务并打开工作台，访问令牌随启动流程传入。关闭窗口或终端前应确认保存状态，并按当前启动方式退出；各平台安装包需分别完成构建与环境验证。

首次使用时，先在“班级与设置”里建立或选择学期和班级，再添加学生；随后用少量合成数据检查导入和统计，最后按需配置 AI 服务。未配置模型时，基础数据功能和 AI 分析的可用性应分别判断。

\subsection{建立班级与导入考试}

\begin{enumerate}
\item 进入“班级与设置”，核对当前学期和班级；在“学生管理”中检查学生名单及学号。
\item 进入“成绩管理”，创建考试并填写名称、满分、阈值等设置。
\item 选择系统支持的成绩表，确认姓名、学号、班级和分数字段的对应关系。需要年级对比时，提供相应的年级数据。
\item 对空缺分数、缺考、重复学号或无法识别的列进行核查；缺考与有效零分必须区分。
\item 保存后查看人数、平均分和分布，抽查部分学生的记录；存在差异时先修正数据，再运行分析。
\end{enumerate}

\subsection{日常记录与错题管理}

在“默写成绩”“背诵成绩”“写作成绩”“日常作业”中建立对应的轮次或任务，录入学生表现并核对。原卷附件和错题记录在“原卷与错题”中管理，填写来源、题型、考点和必要说明。

对话附件、正式考试和正式评价采用不同的数据管理流程。把资料添加到 AI 对话后，应根据实际用途确认是否还需要导入为正式业务记录，避免对话分析与教学记录不一致。

\subsection{发起考试分析}

\begin{enumerate}
\item 进入 TeachMate，选择考试分析功能，核对学期、班级和考试上下文。
\item 如需引用原卷或其他资料，先检查资料内容和选中范围。
\item 输入需求，例如：“请分析当前考试，说明有数据支持的主要薄弱点，并给出讲评建议；材料不足的部分请明确说明。”
\item 查看运行进度。遇到缺少模型配置、材料不足或任务失败时，根据提示处理，不把空结果当作正常完成。
\item 阅读报告中的发现和建议，查看相应证据，核对统计与题目来源。
\item 修改不适合本班的内容，再决定是否采纳或确认。
\end{enumerate}

\subsection{生成学生画像与复习计划}

单人诊断时，教师先选定学生和考试，再发起请求。批量诊断时，先查看待分析名单和执行方案，确认范围及预计消耗后开始；运行过程中可按界面提供的操作停止或重试任务。

阅读画像时，先看优势，再看薄弱点及其依据。对于“可能存在审题问题”这类推断，应结合原题、作答和课堂观察核查。教师确认后的反馈，可以用于个别辅导或家校沟通。

生成复习计划时，可以说明可用课时、关注的知识点及练习负担要求，例如：“根据这次分析，安排一周复习，优先处理最需要补充的两个知识点，并说明每项任务的依据。”实施后由教师记录练习或复测表现，继续比较变化。

\subsection{使用成长森林}

从“学生管理”进入“成长森林”，查看当前班级和学期的成长状态。教师查看学生的奖励来源、阶段和学业表现分支，结合实际教学过程解释其含义。

记录学习活动时填写必要的依据；发现误录时采用更正或撤销流程，不用无来源的重复加分代替修正。展示时强调学生自身的进步，避免把成长点数又变成单一排名。

\subsection{备份、异常处理与退出}

定期保存完整备份；更换设备或恢复数据前，确认数据库、附件和校验信息齐全。普通成绩导出适合查看或交流，但不能替代包含全部业务数据的备份。

\begin{table}[htbp]
\centering
\small
\begin{tabularx}{\textwidth}{@{}L{5.4cm}Y@{}}
\toprule
\textbf{遇到的情况} & \textbf{建议处理} \\
\midrule
程序无法启动 & 查看启动提示，检查依赖、端口占用和运行时准备情况 \\
导入人数或统计异常 & 检查学期、班级、重复学号、满分及缺考状态 \\
AI 按钮不可用或分析失败 & 查看模型配置、网络、预算和当前任务的资料范围 \\
报告缺少逐题分析 & 补充逐题成绩和题目材料，避免从总分推断详细错因 \\
保存发生冲突 & 重新加载最新状态，确认后再次修改 \\
成长值与预期不同 & 查看事件来源、封顶规则、归档和更正记录 \\
\bottomrule
\end{tabularx}
\caption{常见情况与建议处理}
\end{table}

% ===================== 七 =====================
\section{工程验证、应用评估与局限}

\subsection{已有验证记录}

依据 2026 年 9 月 26 日的项目检查记录，相关验证结果如下。验证范围为软件功能与运行流程，教学应用效果另行评估。

\begin{table}[htbp]
\centering
\small
\begin{tabularx}{\textwidth}{@{}L{4.2cm}L{3.2cm}Y@{}}
\toprule
\textbf{检查范围} & \textbf{已记录结果} & \textbf{能支持的结论} \\
\midrule
本地前端逻辑测试 & 152 项通过 & 已覆盖的前端逻辑场景通过检查 \\
成长树与学生快照导入相关测试 & 56 项通过 & 已覆盖的奖励恢复、重复同步、导入字段处理等场景通过检查 \\
本地浏览器端到端测试 & 23 项通过，1 项按设计跳过 & 已覆盖的浏览器流程通过检查；不代表所有平台和模型均已验证 \\
\bottomrule
\end{tabularx}
\caption{已有工程验证记录}
\end{table}

本次修复针对考试或出勤恢复后的成长奖励同步，以及学生快照未提供性别字段时保留既有值等问题。公开仓库已补入相关回归测试，但只包含完整本地测试集合的一部分。

这些检查没有重新构建全部平台的安装包，也没有在本次检查中调用真实付费模型服务。证据引用检查、页面流程检查与真实模型的回答质量，属于不同层面的验证。

\subsection{教学试用评估计划}

后续计划在取得适当授权后，从单教师、少量班级的试用开始：使用相同类型的成绩整理和讲评任务，对比教师原有流程与工作台流程，记录任务完成时间、导入纠错次数、AI 建议被修改或采纳的情况，以及复测记录的完整程度。

学生层面重点观察他们是否理解成长反馈、是否愿意完成订正和复习、任务负担是否合适。评价采用实际记录与访谈结合的方式，不单靠一次主观评价。

试用前会明确任务范围、评价口径及记录方式，并保留失败案例与修改记录。当前尚无系统性教学试验结果；任务效率、学习表现和用户体验等指标，将在实际试用后报告。

\subsection{当前局限}

当前版本主要服务单教师本地使用，跨设备同步、多人协同和独立学生端仍需建设。数据质量直接影响分析质量，复杂附件识别和错因判断需要教师核验。AI 服务依赖配置、网络和供应方状态，模型调用存在成本。成长奖励规则目前属于产品初值，仍需检验其长期激励效果和使用公平性。

此外，项目在成熟度、基础设施与学科覆盖上仍有明显不足：

\begin{itemize}
\item \textbf{尚未成熟，仍处小范围试用。} 当前版本只在有限场景中做过验证，整体并不完善。许多边界情况、交互细节与真实教学的匹配度，需要更多教师在实际使用中暴露问题并反馈，才能逐步修复和完善；这一过程本身也是项目下一步的重点。
\item \textbf{缺少统一的数据库与服务器。} 目前各用户的数据独立保存在本地，没有集中式的数据库或服务端来统一管理所有用户的数据。这有利于数据控制，但也意味着跨用户汇总、统一升级和集中支持尚不具备基础，规模化运营能力仍需建设。
\item \textbf{仅适配单一学科。} 当前只服务初中英语。其他学科在数据结构、知识条目和评价规则上存在差异，扩展到更多学科需要满足不同教师的实际工作方式与字段要求，暂不在当前版本范围内。
\end{itemize}

% ===================== 八 =====================
\section{设计特点与应用价值}

\subsection{教学数据与教学行动的衔接}

系统通过统一的教学数据关联，把平时记录、考试结果、原卷与学生信息用于同一任务；智能体在已选范围内读取数据，形成面向讲评、辅导和复习的建议。其应用价值主要通过任务完成效率、建议可用性和后续执行情况来评估。

\subsection{可追溯的智能体分析}

系统以任务范围约束数据访问，以证据编号关联分析依据，以结构化报告组织结果，使数据来源、分析发现和建议之间具有可检查的联系。教师据此审查结论的适用性，并保留最终的采纳与修改权限。

\subsection{学习过程与学业表现的分层反馈}

成长森林分别呈现学习投入和有证据支持的学业表现，以可比进步和持续参与丰富反馈内容。事件来源与更正记录让奖励可解释，为教师开展过程性评价提供辅助依据。

\subsection{面向单教师场景的渐进式部署}

项目以单教师本地使用为起点，围绕实际任务逐步验证功能与部署要求。初中英语是当前的应用范围；其他学科的适配需要进一步调整数据字段、知识条目和评价规则，并开展相应验证。

上述设计的实际价值，将通过教师完成任务的过程、建议修改记录和后续教学试用进一步评估。

\subsection{低门槛、可扩展与本地可控的设计取向}

项目在易用性、可扩展性和数据控制之间做了明确取舍，这几点是当前版本相对同类方案的主要优势。

\textbf{开箱即用的低门槛。} 智能体运行所需的基础指令与工具调用已在软件内预置，模型服务通常只需填入 API Key 即可启用；TeachMate 各插件内置任务提示词，教师无需手工编写分析要求，降低了非技术背景教师使用 AI 的障碍。

\textbf{预留的数据接口。} 系统从设计上保留了与外部数据源和上层智能体的对接空间：既可接入学校的数据接口，把成绩等业务数据统一纳入工作台分析，也可把教学数据对接到 WorkBuddy 等其他智能体，用于更大范围的编排与协作，避免工具被局限在单一应用内。

\textbf{本地可控与数据安全。} 教学数据完全保存在本地，教师对数据存储、备份与导出拥有直接控制权，日常记录与本地统计无需依赖云端即可运行。

\textbf{保留用户自由。} 软件不把能力锁死在官方版本内：用户可以自行安装插件，也可以通过智能体对软件本身进行改造与增强，使工具随实际使用场景持续演进。

\begin{figure}[htbp]
\centering
\begin{tikzpicture}[font=\footnotesize]
\node[card] (c1) at (-5.4,0) {\textbf{开箱即用的低门槛}\\[4pt]{\scriptsize 基础指令与插件提示词内置}\\[1pt]{\scriptsize 填入 API Key 即可使用}};
\node[card] (c2) at (-1.8,0) {\textbf{预留的数据接口}\\[4pt]{\scriptsize 可接入学校数据接口}\\[1pt]{\scriptsize 也可对接其它智能体}};
\node[card] (c3) at (1.8,0) {\textbf{本地可控与安全}\\[4pt]{\scriptsize 教学数据完全保存在本机}\\[1pt]{\scriptsize 备份导出由教师掌握}};
\node[card] (c4) at (5.4,0) {\textbf{保留用户自由}\\[4pt]{\scriptsize 可自行安装插件}\\[1pt]{\scriptsize 可用智能体改造软件}};
\end{tikzpicture}
\caption{当前版本的四个主要设计取向}
\end{figure}

这些取向的目标，是让工具先能被普通教师用起来、接得进学校既有数据、又能随需要生长，而不是一开始就追求功能完备。

% ===================== 九 =====================
\section{应用前景与商业模式}

\subsection{推广路径}

第一阶段开展小范围教师试用，重点记录数据导入错误、操作障碍和报告修改情况，并据此确定功能迭代优先级。

第二阶段依据试用结果，与教研组或学校评估扩大使用的可行性，并形成配套的操作示范、数据准备说明与服务范围。

第三阶段再评估跨设备、学校内协作和更多学科扩展。扩大使用范围之前，应先完善部署、支持和访问控制能力。

\subsection{拟议商业模式}

当前项目处于验证阶段，尚未建立付费用户、收入或采购合同。可探索的模式如下。

\begin{table}[htbp]
\centering
\small
\begin{tabularx}{\textwidth}{@{}L{3.6cm}YY@{}}
\toprule
\textbf{模式} & \textbf{提供的内容} & \textbf{需要先验证的条件} \\
\midrule
个人基础工具与可选服务 & 基础数据管理，以及自愿选择的 AI 服务、安装或培训支持 & 教师使用意愿、服务成本和长期维护能力 \\
学校部署与支持 & 安装配置、培训、备份方案和版本维护 & 校内授权、实际需求及对多人管理能力的要求 \\
教研场景定制 & 导入模板、校本知识条目和学科流程适配 & 定制工作量、内容来源和交付验收范围 \\
\bottomrule
\end{tabularx}
\caption{拟议的商业模式}
\end{table}

AI 费用可由用户自行配置服务承担，也可在未来服务方案中设置明确额度；具体收费、额度和价格尚未确定。项目不以出售教师或学生数据作为商业路径。

公开源码用于支持实现复核与问题反馈。后续若按开源模式推广，还需明确项目自身的许可安排，并保留第三方组件声明。

\subsection{成本与持续维护}

主要成本包括模型调用、功能开发、跨平台适配、用户支持和教学内容维护。本地部署可减少部分服务端建设需求，但仍需考虑设备、安装、备份和持续维护的投入。

持续维护应优先保障教师已有数据和日常工作：修复影响数据正确性的问题，稳定导入与备份，改进报告的可解释性，再逐步增加新功能。

% ===================== 十 =====================
\section{后续计划与结语}

近期重点是完善参赛演示和源码运行说明，使用合成数据展示完整操作，并补充真实模型运行的可核查记录。随后开展小范围教师试用，检查功能是否真正融入教学工作。

中期计划围绕“分析后如何练习、练习后如何复测”补齐补偿训练闭环，增加题目审核、训练实施和复测对照能力。音标交互手册、听力音频、学生端和学校协同作为独立规划，在确认需求和可行性后逐项推进。

English Workbench 面向初中英语教学中的数据整合、差异化教学和过程反馈需求，构建由本地数据管理、证据支持的智能分析、教师审核及成长记录组成的教学辅助流程。项目已形成可运行的 Beta 原型，后续将以实际教学任务为基础评估使用效果，持续完善从数据记录到教学行动的衔接，为教师开展有依据的教学调整和学生成长反馈提供支持。

% ===================== 材料与实现依据 =====================
\section*{材料与实现依据}
\addcontentsline{toc}{section}{材料与实现依据}

\begin{enumerate}
\item 项目发起人提供的《项目背景与需求分析\_英语教学工作台》：用于教学场景、用户需求和规划方向。
\item 用户提供的 iCAN 复赛提交页面截图：应用方案为 PDF、不超过 20 页；演示视频为 MP4、不超过 5 分钟；单个上传文件不超过 200 MB；软件作品提供可运行链接或应用源码。
\item English Workbench 源码，提交 \texttt{3ad4fe3546e0608c339cd3acf2c54a6213daf088}：启动流程、业务模块、智能体工具、隐私处理与成长规则以实现为依据。仓库：\url{https://github.com/JiajunTang2006/English-workbench}。
\item 2026 年 9 月 26 日项目代码检查及修复记录：用于本文工程验证结果。测试结果不作为教学成效证据。
\end{enumerate}

\end{document}
